\documentclass[11pt,a4paper]{article}
\usepackage[T1]{fontenc}
\usepackage[utf8]{inputenc}
\usepackage[provide=*,main=french]{babel}
\usepackage{lmodern,microtype,geometry,xcolor,hyperref}
\geometry{top=2.5cm,bottom=2.5cm,left=2.7cm,right=2.7cm}
\definecolor{deepblue}{RGB}{21,101,192}
\definecolor{deeporange}{RGB}{216,67,21}
\hypersetup{colorlinks=true,urlcolor=deeporange}
\setlength{\parindent}{0pt}
\setlength{\parskip}{0.75em}
\pagestyle{empty}

\begin{document}
\begin{center}
\vspace*{1.0cm}
{\Huge\bfseries Dynamiques sphériques de Transformers\par}
\vspace{0.45cm}
{\LARGE Dossier théorique et numérique complet\par}
\vspace{0.3cm}
{\large Équilibres, polygones, cycles, MLP, chaos, multi-têtes,\
grande dimension, champ moyen, entraînement et poids variables\par}
\vspace{0.8cm}
{\large Version du 16 août 2026\par}
\end{center}

\vspace{0.55cm}
{\color{deepblue}\hrule}
\vspace{0.55cm}

\textbf{Volume I --- Attention sans MLP, pages 2 à 17.}
Condition nécessaire et suffisante de tous les équilibres, couche finie et temps
continu, réduction spectrale--Gram dans le cas $Q^{\mathsf T}K=V=V^{\mathsf T}$,
formes polygonales, comparaison précise avec Altafini et cycle interne stable produit
par l'attention seule lorsque score et valeur sont déliés.

\textbf{Volume II --- Bloc sériel attention puis MLP, pages 18 à 37.}
Équations exactes de tous les points fixes et cycles, taxonomie des quatre familles,
audit jusqu'à des couches 4 096 fois plus faibles, cycles continus, chaos et
hyperchaos, cartes multi-têtes, limite à beaucoup de tokens, autres noyaux,
observables, entraînement du MLP, première étape de Muon et poids OU lents.

\vspace{0.35cm}
{\color{deepblue}\hrule}
\vspace{0.55cm}

\textbf{Échelle de validation.}
Le dossier s'appuie sur 214 528 modèles tirés indépendamment, 2 844 160 trajectoires
initiales, plus de $10^{11}$ mises à jour token--couche dans l'audit de petite couche,
des intégrations continues indépendantes, des spectres complets de perturbations et
des contrôles frais allant jusqu'à 16 tokens.

\textbf{Règle de lecture.}
Les équations encadrées sont des identités ou conditions exactes. Les tableaux
numériques sont des observations reproductibles sur les protocoles indiqués. Les
idées MaxCut, répliques, résolution universelle de problèmes difficiles et conception
arbitraire d'équilibres sont conservées comme pistes et ne sont pas présentées comme
des théorèmes.

\vfill
\begin{center}
{\small Soumission anonyme --- sources, scripts et agrégats machine inclus dans le
dossier de travail.}
\end{center}
\end{document}
